% ----- Consignes exo 54 ----- %
\begin{td-exo}[Un algorithme de branchement pour le \textsc{Stable Maximum}]\,\\ % 54 
	Proposez un algorithme de branchement pour le problème du \textsc{Stable Maximum}.
\end{td-exo}

% ----- Solutions exo 54 ----- %
\iftoggle{showsolutions}{ 
	\begin{td-sol}[]\ % 54
		On cherche \(\alpha(G)\) le nombre de stabilité du graphe \(G\) (la taille du plus grand stable dans \(G\)).
		On cherche alors
		\begin{equation*}
			\alpha(G) = \max \left\{
				\alpha(V \setminus \{v\}),
				\alpha(V \setminus N(v)) + 1
			\right\}
		\end{equation*}
		En effet, soit \(v\) est dans le stable maximum, alors on ne peut pas prendre les voisins de \(v\) et on a un gain de \(1\) (pour \(v\)) ou alors \(v\) n'est pas dans le stable maximum et on peut prendre les voisins de \(v\).

		On rappelle que pour \(\Delta \leq 2\) on peut trouver un stable maximum en temps linéaire.
		Pour \(\Delta \geq 3\) le problème est NP-complet.
		Notre branchement se fait alors comme suit:
		\begin{equation*}
			\begin{cases}
				T(n) 
				&\leq T(n - 1) + T(n - \Delta - 1)\\
				&\leq T(n - 1) + T(n - 4) 
			\end{cases}
		\end{equation*}
		En se référant à la formule de branchement (dans le cours), on trouve que \(T(n) \leq O(1.3803^n)\).
	\end{td-sol}
}{}
